\section{Limitations}
\label{sec:limitations}

\paragraph{One corpus, one domain.} All measurements use a single English
Wikipedia snapshot chunked at $\sim$100 words. The failure-cause
distribution that motivates the cascade --- and therefore the choice of
its stages --- will differ on enterprise corpora, other languages, and
other chunkings; the diagnosis method transfers, the numbers do not.

\paragraph{The abstention detector is heuristic.} A refusal is recognised
by phrase matching plus a length threshold. Both readers we tested
abstain in easily recognised English formulas; a reader fine-tuned never
to refuse, or prompted in another language, would starve the trigger, and
the cascade would degrade to its first stage. The losslessness property
is likewise specific to span-style Exact Match: under token-level F1 it
is only approximate, and under long-form or generative evaluation it
would not hold.

\paragraph{Two implementation deviations are inherited by the
comparison.} Our CRAG is offline and corpus-only (no web search, no
knowledge refinement) and our Adaptive-RAG classifies with the generator
rather than a trained classifier. Both deviations keep the comparison
controlled --- no row gets an extra model or an extra knowledge source ---
but they also mean our rows bound those methods' training-free, closed-corpus
forms, not the published systems.

\paragraph{Structure-augmented methods are absent.} GraphRAG-family
methods were excluded because their corpus-preprocessing cost is of a
different class at 21M passages. On corpora where that preprocessing is
affordable, the comparison should be re-run with them as rows.

\paragraph{Model scale.} Both readers are 7--8B instruct models. Larger
readers abstain differently --- typically less --- and shift the
closed-book floor upward; the cascade's escalation margin at larger
scales is unmeasured.

\section{Conclusion}
\label{sec:conclusion}

We presented a controlled comparison of thirteen retrieval-augmented
generation architectures --- one corpus, one retriever, one generator,
paired statistics, budget-matched controls, and a measured corpus ceiling
--- together with a per-question diagnosis of why retrieval fails, and an
architecture assembled from that diagnosis. CascadeRAG retrieves with a
fusion of dense, lexical, and cross-encoder rankings, and escalates a
question to a more expensive strategy only when the reader itself, having
seen the evidence, declines to answer --- a posterior, and under Exact
Match provably lossless, trigger. Across \ResultPairs{} paired
comparisons it wins \ResultWins{}, loses \ResultLosses{}, and ties
\ResultTies{}, at 1.16--1.58 generator calls per question, and the result
replicates under a second reader. Alongside the positive result, the
study documents two useful negative ones: neighbour-chunk expansion adds
nothing once context size is controlled, and answer-token confidence adds
nothing over the abstention it accompanies.

The broader claim is methodological. RAG architectures are usually
compared as point scores on incompatible stacks; held to the standard of
paired tests, matched budgets, abstention accounting, and a corpus
ceiling, most published-style gains shrink and one reverses. The
architecture that survives that standard is not the cleverest one
measured but the one that spends its calls only where a free, reliable
signal says they are needed. We release the benchmark, the per-question
logs, and the manuscript pipeline so that both the standard and the
architecture can be applied to other stacks as they are.
